% Options for packages loaded elsewhere
% Options for packages loaded elsewhere
\PassOptionsToPackage{unicode}{hyperref}
\PassOptionsToPackage{hyphens}{url}
\PassOptionsToPackage{dvipsnames,svgnames,x11names}{xcolor}
%
\documentclass[
  english,
  11pt,
]{article}
\usepackage{xcolor}
\usepackage[margin=25mm]{geometry}
\usepackage{amsmath,amssymb}
\setcounter{secnumdepth}{-\maxdimen} % remove section numbering
\usepackage{iftex}
\ifPDFTeX
  \usepackage[T1]{fontenc}
  \usepackage[utf8]{inputenc}
  \usepackage{textcomp} % provide euro and other symbols
\else % if luatex or xetex
  \usepackage{unicode-math} % this also loads fontspec
  \defaultfontfeatures{Scale=MatchLowercase}
  \defaultfontfeatures[\rmfamily]{Ligatures=TeX,Scale=1}
\fi
\usepackage{lmodern}
\ifPDFTeX\else
  % xetex/luatex font selection
\fi
% Use upquote if available, for straight quotes in verbatim environments
\IfFileExists{upquote.sty}{\usepackage{upquote}}{}
\IfFileExists{microtype.sty}{% use microtype if available
  \usepackage[]{microtype}
  \UseMicrotypeSet[protrusion]{basicmath} % disable protrusion for tt fonts
}{}
\makeatletter
\@ifundefined{KOMAClassName}{% if non-KOMA class
  \IfFileExists{parskip.sty}{%
    \usepackage{parskip}
  }{% else
    \setlength{\parindent}{0pt}
    \setlength{\parskip}{6pt plus 2pt minus 1pt}}
}{% if KOMA class
  \KOMAoptions{parskip=half}}
\makeatother
% Make \paragraph and \subparagraph free-standing
\makeatletter
\ifx\paragraph\undefined\else
  \let\oldparagraph\paragraph
  \renewcommand{\paragraph}{
    \@ifstar
      \xxxParagraphStar
      \xxxParagraphNoStar
  }
  \newcommand{\xxxParagraphStar}[1]{\oldparagraph*{#1}\mbox{}}
  \newcommand{\xxxParagraphNoStar}[1]{\oldparagraph{#1}\mbox{}}
\fi
\ifx\subparagraph\undefined\else
  \let\oldsubparagraph\subparagraph
  \renewcommand{\subparagraph}{
    \@ifstar
      \xxxSubParagraphStar
      \xxxSubParagraphNoStar
  }
  \newcommand{\xxxSubParagraphStar}[1]{\oldsubparagraph*{#1}\mbox{}}
  \newcommand{\xxxSubParagraphNoStar}[1]{\oldsubparagraph{#1}\mbox{}}
\fi
\makeatother


\usepackage{longtable,booktabs,array}
\newcounter{none} % for unnumbered tables
\usepackage{calc} % for calculating minipage widths
% Correct order of tables after \paragraph or \subparagraph
\usepackage{etoolbox}
\makeatletter
\patchcmd\longtable{\par}{\if@noskipsec\mbox{}\fi\par}{}{}
\makeatother
% Allow footnotes in longtable head/foot
\IfFileExists{footnotehyper.sty}{\usepackage{footnotehyper}}{\usepackage{footnote}}
\makesavenoteenv{longtable}
\usepackage{graphicx}
\makeatletter
\newsavebox\pandoc@box
\newcommand*\pandocbounded[1]{% scales image to fit in text height/width
  \sbox\pandoc@box{#1}%
  \Gscale@div\@tempa{\textheight}{\dimexpr\ht\pandoc@box+\dp\pandoc@box\relax}%
  \Gscale@div\@tempb{\linewidth}{\wd\pandoc@box}%
  \ifdim\@tempb\p@<\@tempa\p@\let\@tempa\@tempb\fi% select the smaller of both
  \ifdim\@tempa\p@<\p@\scalebox{\@tempa}{\usebox\pandoc@box}%
  \else\usebox{\pandoc@box}%
  \fi%
}
% Set default figure placement to htbp
\def\fps@figure{htbp}
\makeatother



\ifLuaTeX
\usepackage[bidi=basic,shorthands=off]{babel}
\else
\usepackage[bidi=default,shorthands=off]{babel}
\fi
\ifLuaTeX
  \usepackage{selnolig} % disable illegal ligatures
\fi


\setlength{\emergencystretch}{3em} % prevent overfull lines

\providecommand{\tightlist}{%
  \setlength{\itemsep}{0pt}\setlength{\parskip}{0pt}}



 


\makeatletter
\@ifpackageloaded{caption}{}{\usepackage{caption}}
\AtBeginDocument{%
\ifdefined\contentsname
  \renewcommand*\contentsname{Table of contents}
\else
  \newcommand\contentsname{Table of contents}
\fi
\ifdefined\listfigurename
  \renewcommand*\listfigurename{List of Figures}
\else
  \newcommand\listfigurename{List of Figures}
\fi
\ifdefined\listtablename
  \renewcommand*\listtablename{List of Tables}
\else
  \newcommand\listtablename{List of Tables}
\fi
\ifdefined\figurename
  \renewcommand*\figurename{Figure}
\else
  \newcommand\figurename{Figure}
\fi
\ifdefined\tablename
  \renewcommand*\tablename{Table}
\else
  \newcommand\tablename{Table}
\fi
}
\@ifpackageloaded{float}{}{\usepackage{float}}
\floatstyle{ruled}
\@ifundefined{c@chapter}{\newfloat{codelisting}{h}{lop}}{\newfloat{codelisting}{h}{lop}[chapter]}
\floatname{codelisting}{Listing}
\newcommand*\listoflistings{\listof{codelisting}{List of Listings}}
\makeatother
\makeatletter
\makeatother
\makeatletter
\@ifpackageloaded{caption}{}{\usepackage{caption}}
\@ifpackageloaded{subcaption}{}{\usepackage{subcaption}}
\makeatother
\usepackage{bookmark}
\IfFileExists{xurl.sty}{\usepackage{xurl}}{} % add URL line breaks if available
\urlstyle{same}
% fallback for those not using the hyperref driver hyperxmp:
\makeatletter
\@ifundefined{xmpquote}{\newcommand{\xmpquote}[1]{#1}}{}
\makeatother
\hypersetup{
  pdftitle={The missing push: what a residual means},
  pdflang={en},
  colorlinks=true,
  linkcolor={blue},
  filecolor={Maroon},
  citecolor={Blue},
  urlcolor={Blue},
  pdfcreator={LaTeX via pandoc}}


\title{The missing push: what a residual means}
\usepackage{etoolbox}
\makeatletter
\providecommand{\subtitle}[1]{% add subtitle to \maketitle
  \apptocmd{\@title}{\par {\large #1 \par}}{}{}
}
\makeatother
\subtitle{A force-balance question behind the whole construction}
\author{}
\date{2026-09-12}
\begin{document}
\maketitle


\href{../paper-map/article.html}{Start with the paper map} · Main route,
article 4 of 9

We can draw a vortex that gets narrower and faster. We can even arrange
its geometry so that the energy stays bounded. But does that motion
actually satisfy the fluid equations?

To answer, we calculate the push it would require. The result is called
a \textbf{residual}. It tells us what is left in the force balance after
the chosen velocity and pressure have been inserted.

This is the central question connecting the background vortex to the
later oscillations. We will begin with a cart, then translate the same
accounting into fluid motion.

\subsection{One kilogram, three newtons, and a missing
newton}\label{one-kilogram-three-newtons-and-a-missing-newton}

Newton's law says that force equals mass times acceleration. A
one-kilogram cart accelerating at three metres per second squared needs
a total force of three newtons.

Suppose the forces we have already included add up to two newtons. One
newton is missing:

\[
\text{missing push}=3-2=1.
\]

The cart's proposed motion might be perfectly smooth. That does not
explain the missing newton. We must supply it or change the motion.

\begin{figure}

\centering{

\pandocbounded{\includegraphics[keepaspectratio,alt={Bars show a required force of three, an accounted-for force of two, and a missing force of one.}]{figures/reader.png}}

}

\caption{\label{fig-balance}A one-kilogram cart needs three newtons for
the stated acceleration. Two are accounted for and one remains. The bar
diagram illustrates a residual as an accounting difference.}

\end{figure}%

\subsection{The fluid has several contributions to account
for}\label{the-fluid-has-several-contributions-to-account-for}

A small parcel of fluid can change its velocity for several reasons.

Pressure differences push it. Viscosity exchanges momentum with
neighbouring fluid moving differently. An external force may act. The
parcel also moves into new places, where the surrounding velocity may
differ from where it began.

That last point is easy to overlook. A car entering a tighter bend
changes its motion partly because it has reached a different part of the
road. Following a moving parcel is different from watching one fixed
point in space.

The Navier--Stokes equation combines these effects. For our purpose, we
can read it as a sentence:

\begin{quote}
The acceleration of a fluid parcel must match the pressure, viscous, and
external contributions.
\end{quote}

The paper's momentum residual is the \textbf{external force per unit
mass required by the chosen velocity and pressure}. If a particular
external force has already been specified, subtracting it gives the
remaining equation mismatch.

A nonzero residual is therefore not automatically a mistake. A forced
solution can have a nonzero required force. What matters here is whether
that force has the smoothness and support required by the theorem.

\subsection{A constant speed can still require
acceleration}\label{a-constant-speed-can-still-require-acceleration}

Imagine moving around a circle at constant speed. Your direction keeps
changing, so your velocity changes. An inward acceleration is required
even though the speedometer stays fixed.

This is why a rotating flow needs a pressure balance in the radial
direction. A calculation that treats the circular velocity as an
ordinary scalar speed and ignores the changing direction misses a real
force.

The \href{../d6_axis_pressure/article.html}{pressure article} explained
that part of the balance. But satisfying the inward circular balance
does not automatically satisfy the balances for changes in speed, axial
motion, and viscosity.

A single successful cancellation is one entry in the account, not the
whole account.

\subsection{Why simply defining a force does not finish the
paper}\label{why-simply-defining-a-force-does-not-finish-the-paper}

Given a candidate flow, one can formally calculate the force that makes
its equation hold. That sounds like an easy way to construct almost any
motion.

The important restriction is on the resulting force. The paper requires
a smooth force that remains well behaved through the time when the
velocity becomes unbounded.

The background vortex alone leaves a troublesome residual in the joining
annulus. If that leftover becomes singular, merely naming it ``the
force'' does not meet the target.

The paper therefore changes the velocity itself. Its waves create
additional internal momentum transport, chosen to cancel the troublesome
leading contribution. Later corrections deal with the rest.
\href{https://cdn.openai.com/pdf/32d9f210-8b73-45e0-91bc-82a30aef8a9a/navier-stokes.pdf}{Paper,
Section 2.2 and Sections 3.2--3.4}

\subsection{Two examples that make the distinction
concrete}\label{two-examples-that-make-the-distinction-concrete}

Our supporting calculation compares an explicit outer flow from the
paper with a chosen smooth vortex-shaped example.

The outer flow has a special property: its circulating velocity follows
a diffusion equation, and its pressure provides the needed circular
acceleration. The relevant momentum residual cancels. The source calls
this the \emph{heat exterior} because the smoothing equation is related
to the heat equation. This label is about a mathematical form, not a
calculation of the fluid's temperature.

The second example has a prescribed shrinking shape and a growing peak.
Those features make it useful for explaining concentration. They do not
automatically make its force balance work. Its remaining residual
exposes that difference.

These calculations show why a convincing animation is insufficient
evidence that a field solves the desired equation. We must calculate
what force the animation's motion requires.

\subsection{Why large terms can still leave a smooth
result}\label{why-large-terms-can-still-leave-a-smooth-result}

Suppose one contribution is 1,000,001 and another subtracts 1,000,000.
The result is 1. Individual terms can be large while their combination
is small.

The paper relies on carefully organized cancellation of large
contributions. The challenge is to keep that cancellation valid as the
core shrinks and to control how the remaining force changes in space and
time.

A few displayed small numbers cannot establish all of that. But the
elementary subtraction explains why a growing velocity and a smooth
force are not logically incompatible.

\subsection{Where this fits in the bigger
picture}\label{where-this-fits-in-the-bigger-picture}

\textbf{The residual identifies the repair target.} It turns a proposed
vortex shape into a concrete question about the transport that is
missing.

The next blocks explain how fine structure helps meet local constraints
and how oscillating motion supplies averaged transport. Continue with
\href{../convex-integration/article.html}{convex integration, from
slopes to fluid motion}.

The \href{technical-article.html}{technical article},
\href{../../code/d3_residual/study.ipynb}{notebook}, and
\href{source-notes.md}{source notes} contain the actual momentum terms
and the two computational examples.




\end{document}
